\subsection{Case 1 : $D_i$ is exogenous to $N, Y$}

Figure \ref{fig:simcase1} \rev{and Tables~\ref{tab:AppSimExoCase1}--\ref{tab:AppSimEndCase1} show} that even in case the exogeneity of deductible is assured both for the frequency and severity, usage of either FEM2 or SEM2/SEM3 is less efficient in terms of the estimation performance so that the RMSE are larger than the ones from FEM1 and SEM1, respectively. As FEM2 depends on $p(d;\bm{x})$ from the severity models, FEM2G suffers less than FEM2. \rev{The joint models do not find spurious dependence: the estimates of $\rho_Y$ in SDM1--3 and of $\rho_N$ in FDM1--2 are centered at zero, their regression coefficients are approximately unbiased, and the price paid for estimating an unnecessary dependence parameter is small, with RMSEs close to those of the corresponding marginal models (Table~\ref{tab:AppSimEndCase1}). The ordinal deductible model itself is estimated with negligible bias in all cases (Table~\ref{tab:AppSimDed}).} Table \ref{tab:maecase1} shows that the predictions for the censored and truncated reporting scenarios exhibit comparable levels of estimation error \rev{under all models} when the deductible is exogenous\rev{; the premium errors of the joint models are only slightly larger than those of the marginal models}.

\begin{figure}[h!]
    \centering
    \includegraphics[width=\linewidth]{simcase1.pdf}
    \caption{\rev{Estimates of the severity (upper panels) and frequency (lower panels) parameters in Case 1; the dashed line is the true value and the shaded boxes are the joint models}}
    \label{fig:simcase1}
\end{figure}

@@tab_mae_case1.tex@@

\FloatBarrier

\subsection{Case 2 : $D_i$ is exogenous to $Y$ and $(D, N)$ is from a Gaussian copula with $\rho_N=-0.5$}
Figure \ref{fig:simcase2} shows that even in case the exogeneity of deductible is assured for the severity, usage of  SEM2/SEM3 is less efficient in terms of the estimation performance so that the RMSE are larger than the ones from SEM1.  It is further observed that usage of FEM2 suffers from biases in the estimated parameters as it fails to capture possible endogeneity of the deductible with incomplete data\rev{: the intercept of $\bm{\beta}_N$ is overestimated by about $0.18$ and $r$ by about $0.5$ (Table~\ref{tab:AppSimExoCase2}), i.e., the thinning of the claims by the endogenous deductible is wrongly attributed to the mean and the dispersion of $N$}. Moreover, FEM2G still shows similar bias and RMSE to those of FEM2, which implies that fixing $p(d;\bm{x})$ gives less improvement when the deductible is not exogenous to the frequency. It is also notable that even with incomplete frequency data, \rev{FDM2} still yields \rev{approximately unbiased} estimates of the parameters \rev{including $\rho_N$ (the biases in Table~\ref{tab:AppSimEndCase2} are within a few Monte Carlo standard errors of zero)} by successfully incorporating the existing endogeneity of the deductible for the frequency component. In Table \ref{tab:maecase2}, \rev{FDM1/FDM2} also show lower prediction error in both censored and truncated frequency\rev{, and the error of the frequency component of the premium is reduced by an order of magnitude (from about $0.38$ to $0.03$ claims per policy-year under censored reporting, and from about $0.19$ to $0.02$ under truncated reporting), with FDM2 providing the same level of prediction as FDM1 from complete data}.

\begin{figure}[h!]
    \centering
    \includegraphics[width=\linewidth]{simcase2.pdf}
    \caption{\rev{Estimates of the severity (upper panels) and frequency (lower panels) parameters in Case 2}}
    \label{fig:simcase2}
\end{figure}

@@tab_mae_case2.tex@@

\FloatBarrier

\subsection{Case 3 : $D_i$ is exogenous to $N$ and $(D, Y)$ is from a Gaussian copula with $\rho_Y=-0.5$}
Figure \ref{fig:simcase3} shows that in case the exogeneity of deductible is assured only for the frequency, it is observed that usage of SEM2/SEM3 suffers from more biases in the estimated parameters as it fails to capture possible endogeneity of the deductible with incomplete data\rev{; the biases are large (about $+0.6$ to $+0.7$ for the intercept of $\bm{\beta}_Y$ and about $+1.4$ for $\alpha$, Table~\ref{tab:AppSimExoCase3}), because the losses that are censored or truncated by the deductible are not a random subset of the losses when the deductible is informative about the loss size}. It is also notable that even with incomplete severity data, \rev{SDM2/SDM3} still yield \rev{approximately unbiased} estimates of the parameters \rev{(Table~\ref{tab:AppSimEndCase3})} by successfully incorporating the existing endogeneity of the deductible for the severity component. \rev{Note that in this case FEM2 and FEM2G are also biased, although $D\perp N$ (Table~\ref{tab:AppSimExoCase3}), because the retention probability $p(d;\bm{x})$ used by the truncated marginal frequency models is computed under the exogeneity assumption, i.e., from the marginal rather than the conditional distribution of $Y$ given $D$ (and, for FEM2, from the biased marginal severity model SEM3); this illustrates the indirect channel through which the endogeneity of the deductible with respect to the severity contaminates the frequency model under truncated reporting, and FDM2, which evaluates $p(d;\bm{x})$ with SDM3, is free of it.} In Table \ref{tab:maecase3}, \rev{SDMs} provide lower prediction errors than SEMs and \rev{SDM2/SDM3} show improved prediction \rev{at a similar level to SDM1; the error of the severity component of the premium drops from above $1{,}300$ under SEM1--3 to between about $100$ and $200$ under SDM1--3, and the frequency premium error of FDM2 is about one fifth of that of FEM2}.

\begin{figure}[h!]
    \centering
    \includegraphics[width=1\linewidth]{simcase3.pdf}
    \caption{\rev{Estimates of the severity (upper panels) and frequency (lower panels) parameters in Case 3}}
    \label{fig:simcase3}
\end{figure}

@@tab_mae_case3.tex@@

\FloatBarrier

\subsection{Case 4 : Both $(D,N),(D, Y)$ from Gaussian copulas with $\rho_N=-0.6,  \rho_Y=-0.4$}

Figure \ref{fig:simcase4} shows that in case the exogeneity of deductible is not assured, it is observed that usage of FEM2 and SEM2/SEM3 severely suffers from biases in the estimated parameters as they fail to capture possible endogeneity of the deductible with incomplete data. \rev{In this case even SEM1, which uses the complete ground-up losses, is biased (Table~\ref{tab:AppSimExoCase4}). The reason is that the units of the severity likelihood are the claims: since $N$ and $D$ are dependent, policy-years with a small deductible contribute more claims, and since $Y$ and $D$ are dependent, their losses are larger, so that the pooled distribution of the ground-up losses over the claims is not the marginal distribution of $Y$ given $\bm{x}$ assumed by SEM1. The joint severity models condition on the deductible and are therefore not affected by this size-biased sampling of the claims.}
It is also notable that even with incomplete  data, \rev{FDM2 and SDM2/SDM3} still yield \rev{approximately unbiased} estimates of the parameters \rev{(Table~\ref{tab:AppSimEndCase4})} by successfully incorporating the existing endogeneity of the deductible for both the frequency and severity components. In Table \ref{tab:maecase4}, one can observe that \rev{FDM/SDMs} show lower prediction error than FEM/SEMs and the gap of MAE between \rev{FDM1/SDM1} and \rev{FDM2/SDM2/SDM3} decreases providing comparable level of error with complete data. \rev{Table~\ref{tab:bylevel4} shows what the biases of the marginal models mean for ratemaking. Since the marginal models ignore the deductible, their premium components are too low for the policies with the lowest deductible and far too high for the policies with high deductibles: under censored reporting the frequency component is underestimated by about $42\%$ at the deductible of $500$ and overestimated by about $20\%$ to $440\%$ at the higher levels by FEM1, and the severity component by SEM1/SEM2 shows errors between $-18\%$ and $+680\%$; under truncated reporting the frequency errors of FEM1/FEM2 are even larger (above $+2{,}000\%$ at the highest level, where only a small fraction of the losses exceeds the deductible). The joint models, which price each deductible level with its own conditional distribution, are within about $3\%$ at every level for the frequency and within about $1\%$ for the severity, under complete and incomplete data alike.}

\begin{figure}[h!]
    \centering
    \includegraphics[width=1\linewidth]{simcase4.pdf}
    \caption{\rev{Estimates of the severity (upper panels) and frequency (lower panels) parameters in Case 4}}
    \label{fig:simcase4}
\end{figure}

@@tab_mae_case4.tex@@

@@tab_bylevel_case4.tex@@

\rev{Table~\ref{tab:simcoverage} uses Cases 1 and 4 to check the standard errors of the sequential estimator. For every replication the stage-wise (naive) standard errors, the stage-wise sandwich standard errors and the Godambe (sandwich) standard errors of Section 2.3 are computed and compared with the Monte Carlo standard deviation of the estimates; the table also reports the empirical coverage of the nominal 95\% confidence intervals. For SDM3 the three standard errors practically coincide, i.e., the uncertainty of the deductible model estimated at the first stage is immaterial for the severity model. For the truncated frequency models, which depend on the severity model through $p(d;\bm{x})$, the naive standard errors are too small: for FEM2 in Case 1 and for FDM2 in Case 4 the Godambe standard errors of the regression coefficients, of $r$ and of $\rho_N$ are larger than the naive ones by roughly 10--20\%, and the coverage of the naive intervals tends to fall below the nominal level while that of the sandwich intervals is close to 95\%. This is the finite-sample counterpart of the cross-derivative blocks $\bm{H}_{32}$ in Section 2.3, and it is the reason why all standard errors in Section~\ref{sec:empirical} are computed from the Godambe information matrix.}

@@tab_coverage.tex@@

\FloatBarrier

\subsection{\rev{Cases 5 and 6 : violation of the conditional independence assumption}}

\rev{Cases 5 and 6 keep the dependence parameters of Case 4 and violate the assumption $N\perp Y\,|\,D$ in the two ways described above; the results are given in Tables~\ref{tab:AppSimExoCase5}--\ref{tab:AppSimEndCase6} and Figures~\ref{fig:simcase5}--\ref{fig:simcase6}. In Case 5, where a common latent propensity drives the deductible, the frequency and the severity, the pairs $(D,N)$ and $(D,Y)$ still follow the assumed copulas, and FDM1 remains approximately unbiased. The severity joint models, however, are fitted to claims that concentrate at the low-deductible end of each atom and show moderate biases: about $-0.04$ for the slopes of $\bm{\beta}_Y$, $-0.1$ to $-0.15$ for $\alpha$ and $-0.02$ to $-0.04$ for $\rho_Y$; FDM2 inherits a small bias through $p(d;\bm{x})$ (about $-0.05$ for $r$). These biases are an order of magnitude smaller than those of the marginal models (for example $+0.6$ to $+0.7$ for the intercept of $\bm{\beta}_Y$ and $+0.6$ to $+0.8$ for $\alpha$ under SEM1--3, and $+0.17$ for the intercept of $\bm{\beta}_N$ and $+0.18$ for $r$ under FEM2), and the prediction errors in Table~\ref{tab:maecase5} remain much lower under the joint models. In Case 6, where in addition the frequency and the severity are directly dependent given the deductible, the size-biased sampling of the claims operates within each deductible level, and the joint severity models become substantially biased even with complete data (about $+0.6$ for the intercept of $\bm{\beta}_Y$ and $+0.5$ for $\alpha$ under SDM1), although still less than the marginal models (about $+1.2$ and $+1.5$ under SEM1); the frequency models are more robust, with FDM1 approximately unbiased and FDM2 showing moderate biases (about $+0.04$ for the intercept of $\bm{\beta}_N$, $-0.07$ for $r$ and $+0.04$ for $\rho_N$). In both cases the estimated dependence parameters stay close to their true values, so that the endogeneity of the deductible is still correctly detected; what suffers is the severity regression. A direct frequency--severity dependence given the deductible would therefore have to be modeled explicitly, which, as discussed in Section~\ref{sec:condind}, is not straightforward under censored or truncated losses and is left for future research.}

\begin{figure}[h!]
    \centering
    \includegraphics[width=1\linewidth]{simcase5.pdf}
    \caption{\rev{Estimates of the severity (upper panels) and frequency (lower panels) parameters in Case 5}}
    \label{fig:simcase5}
\end{figure}

\begin{figure}[h!]
    \centering
    \includegraphics[width=1\linewidth]{simcase6.pdf}
    \caption{\rev{Estimates of the severity (upper panels) and frequency (lower panels) parameters in Case 6}}
    \label{fig:simcase6}
\end{figure}

@@tab_mae_case5.tex@@

@@tab_mae_case6.tex@@
